\chapter{The Weil algebra and its domain}
\label{v4-arith:weil:chapter}

The explicit formula equates arithmetic and spectral evaluations of a test
function. The arithmetic evaluation contains every prime power, a gamma
integral and two polar values. The spectral evaluation sums over the zeros
of the completed zeta function. Applying the formula to convolution products
produces a Hermitian form. Its positivity on the full test algebra is
equivalent to the Riemann hypothesis.

The domain is part of this assertion. Compact support makes each prime sum
finite. Exponential decay permits Gaussian tests and controls the infinite
prime sum. The ordinary Schwartz topology controls the gamma integral but
does not control the critically weighted prime measure. We shall construct
the relevant extensions and the maps between their domains.

Two examples explain why these choices matter. The polar form has rank two
and both signs. The finite-prime form also has a two-dimensional subspace
with both signs. Their contributions enter the complete formula with
specified signs. A positive realization must reproduce that complete form.

Fourier inversion, elementary distribution theory and complex analysis are
the prerequisites. We use Lebesgue measure on the additive real line.
Hermitian forms are linear in their first argument. The convolution
algebras below have no unit unless a unitization is explicitly taken.

\section{Compact tests and Mellin reflection}
\label{v4-arith:weil:compact}

For \(R>0\), let \(\mathscr A_R\) be the space of complex smooth functions
supported in \([-R,R]\), with seminorms
\[
 \|f\|_{R,j}=\sup_{x\in\mathbb R}|f^{(j)}(x)|,\qquad j\ge0.
\]
The compact test space is
\[
 \mathscr A_c=C_c^\infty(\mathbb R)
       =\underset{R\to\infty}{\operatorname{colim}}\mathscr A_R
\]
with its usual strict locally convex inductive-limit topology. A convergent
sequence lies in one \(\mathscr A_R\) and converges there in every derivative
seminorm.

\begin{definition}[The compact convolution algebra]
\label{v4-arith:weil:def:compact}
For \(f,g\in\mathscr A_c\), set
\begin{equation}
 (f*g)(x)=\int_{\mathbb R}f(y)g(x-y)\,dy,\qquad
 f^*(x)=\overline{f(-x)}.
 \label{v4-arith:weil:eq:convolution}
\end{equation}
The Fourier transform and its real-axis inverse are
\begin{equation}
 h_f(z)=\int_{\mathbb R}f(x)e^{izx}\,dx,\qquad
 f(x)=\frac1{2\pi}\int_{\mathbb R}h_f(t)e^{-itx}\,dt.
 \label{v4-arith:weil:eq:fourier}
\end{equation}
For an entire function \(h\), write
\(h^\sharp(z)=\overline{h(\bar z)}\).
\end{definition}

\begin{proposition}[Products, support and involution]
\label{v4-arith:weil:prop:algebra}
Convolution and \(f\mapsto f^*\) make \(\mathscr A_c\) a commutative
topological complex \(*\)-algebra. They satisfy
\begin{equation}
 \operatorname{supp}(f*g)\subset
 \operatorname{supp}f+\operatorname{supp}g,\qquad
 h_{f*g}=h_fh_g,\qquad h_{f^*}=h_f^\sharp.
 \label{v4-arith:weil:eq:fourier-algebra}
\end{equation}
For \(f\in\mathscr A_R\), \(h_f\) is entire and obeys
\[
 |h_f(z)|\le C_{R,N}(f)e^{R|\operatorname{Im}z|}
                       (1+|z|)^{-N}\qquad(N\ge0),
\]
where \(C_{R,N}\) is a continuous seminorm. In particular, \(h_f\) and all
its derivatives decrease rapidly on every closed horizontal strip.
\end{proposition}
\begin{proof}
The support assertion follows directly from the integral. Fubini proves
associativity, commutativity and the product identity. Reflection and
conjugation give the involution identity. Differentiation may be performed
under convolution on a fixed compact support. The resulting estimates
prove continuity on the compact test space.

Compact support also permits differentiation of \(h_f\) at every complex
argument. Repeated integration by parts gives
\[
 (iz)^Nh_f(z)=(-1)^N\int f^{(N)}(x)e^{izx}\,dx.
\]
There are no boundary terms. Combine this bound for \(|z|\ge1\) with the
integral bound for \(|z|\le1\). Applying the same argument to \(x^jf(x)\)
proves the assertion about derivatives.
\end{proof}

We give the Fourier image \(h(\mathscr A_c)\) the transported topology.
An entire function in this image cannot vanish on a real interval unless
it is zero. Exact restriction to a bounded frequency interval therefore
does not preserve this Fourier image.

The map \(\epsilon(f)=\int f\) is multiplicative. It does not provide a
convolution unit. Indeed, a unit \(e\in\mathscr A_c\) would satisfy
\(h_e(t)h_f(t)=h_f(t)\) for every \(f\). At each real \(t\), some compact
test has nonzero transform, so \(h_e(t)=1\). This contradicts rapid decay.

The Dirac distribution at zero is the unit in a larger distribution
algebra. Smooth functions \(b_\delta(x)=\delta^{-1}b(x/\delta)\), with
\(\int b=1\), instead form an approximate identity on \(\mathscr A_c\).

The multiplicative coordinate \(u=e^x\) introduces a half-density.
For \(f\in\mathscr A_c\), define
\begin{equation}
 \phi_f(u)=u^{-1/2}f(\log u),\qquad
 \Phi_f(s)=\int_0^\infty\phi_f(u)u^s\,\frac{du}{u}
          =h_f\bigl(-i(s-\tfrac12)\bigr).
 \label{v4-arith:weil:eq:mellin}
\end{equation}
Thus \(\phi_f\) is compactly supported on \(\mathbb R_{>0}\).
Multiplicative convolution uses the measure \(du/u\).

\begin{proposition}[The two reciprocal maps]
\label{v4-arith:weil:prop:reciprocal}
The assignment \(f\mapsto\phi_f\) preserves convolution. Its induced
antilinear involution and its linear reciprocal map are
\begin{equation}
 \phi^*(u)=u^{-1}\overline{\phi(u^{-1})},\qquad
 \phi^\iota(u)=u^{-1}\phi(u^{-1}).
 \label{v4-arith:weil:eq:reciprocal}
\end{equation}
Their Mellin transforms are
\[
 \Phi_{\phi^*}(s)=\overline{\Phi_\phi(1-\bar s)},\qquad
 \Phi_{\phi^\iota}(s)=\Phi_\phi(1-s).
\]
In particular,
\[
 \Phi_{f*g^*}(s)=
       \Phi_f(s)\overline{\Phi_g(1-\bar s)}.
\]
\end{proposition}
\begin{proof}
In multiplicative convolution substitute \(u=e^x\), \(v=e^y\).
The two half-density factors multiply to \(e^{-x/2}\), leaving
\((f*g)(x)\). Substitution \(u\mapsto u^{-1}\) in the Mellin integral
gives both reciprocal identities. The last formula follows from
Fourier multiplication.
\end{proof}

Mellin reflection is linear. A Hermitian pairing also needs complex
conjugation. Equivalently, a form \(q\) linear in its first argument is a
complex-linear functional on \(\mathscr A_c\otimes_{\mathbb C}
\overline{\mathscr A_c}\), where the bar denotes the conjugate vector
space. Replacing this factor by \(\mathscr A_c\) changes the scalar rule:
a bilinear diagonal changes sign under \(f\mapsto if\), whereas a Hermitian
diagonal does not.

\section{Prime powers and exponential weights}
\label{v4-arith:weil:primes}

Let \(\Lambda(n)\) be the von Mangoldt function:
\(\Lambda(p^k)=\log p\) for primes \(p\) and integers \(k\ge1\), and
\(\Lambda(n)=0\) otherwise. The logarithmic prime measures are
\begin{equation}
 \mu_\sigma=\sum_{n\ge2}\Lambda(n)n^{-\sigma}\delta_{\log n},
 \qquad \sigma\in\mathbb R.
 \label{v4-arith:weil:eq:prime-measure}
\end{equation}
Every compact interval meets finitely many atoms. Each \(\mu_\sigma\)
therefore acts continuously on \(\mathscr A_c\).

\begin{lemma}[Elementary prime mass]
\label{v4-arith:weil:lem:prime-mass}
Put \(\Psi(X)=\sum_{n\le X}\Lambda(n)\). For integers \(m\ge1\),
\[
 2m\log2-\log(2m+1)\le\Psi(2m).
\]
Moreover, \(\Psi(X)\le X\log X\) for \(X\ge2\). Consequently,
\(\Psi(X)\ge cX\) for all sufficiently large \(X\), with some \(c>0\).
\end{lemma}
\begin{proof}
The largest coefficient in \((1+1)^{2m}\) is \(\binom{2m}{m}\), so
\(\binom{2m}{m}\ge4^m/(2m+1)\). Its \(p\)-adic valuation is
\[
 v_p\binom{2m}{m}
 =\sum_{k\ge1}\left(
       \left\lfloor\frac{2m}{p^k}\right\rfloor
       -2\left\lfloor\frac m{p^k}\right\rfloor\right).
\]
Each summand is zero or one and vanishes for \(p^k>2m\).
Taking logarithms gives
\(\log\binom{2m}{m}\le\Psi(2m)\). The upper bound follows from
\(\Lambda(n)\le\log n\). Monotonicity, with \(m=\lfloor X/2\rfloor\),
gives the final lower bound.
\end{proof}

\begin{theorem}[The exact temperedness threshold]
\label{v4-arith:weil:thm:prime-threshold}
The positive measure \(\mu_\sigma\) is tempered exactly when
\(\sigma\ge1\). It has finite total mass exactly when \(\sigma>1\).
In particular, the critical measure \(\mu_{1/2}\) is not tempered.
\end{theorem}
\begin{proof}
For \(\sigma\ge1\) and a Schwartz function \(v\),
\[
 |\mu_\sigma(v)|
 \le\sup_x(1+|x|)^3|v(x)|
       \sum_{n\ge2}\frac{\log n}{n(1+\log n)^3}.
\]
The series converges by the integral test. For \(\sigma>1\), the
stronger bound \(\sum(\log n)n^{-\sigma}<\infty\) proves finite mass.
At \(\sigma=1\), partial summation gives
\[
 \sum_{n\le X}\frac{\Lambda(n)}n
 =\frac{\Psi(X)}X+\int_1^X\frac{\Psi(u)}{u^2}\,du.
\]
The lower bound of the lemma makes this quantity diverge.

A positive tempered measure has polynomially bounded mass on
\([-R,R]\). To see this, test against \(v(x/R)\), where \(v\ge0\) is a
fixed compact smooth function equal to one on \([-1,1]\).
Each of its finitely many required Schwartz seminorms grows polynomially
in \(R\).

If \(0\le\sigma<1\), then
\[
 \mu_\sigma([0,R])
 \ge e^{-\sigma R}\Psi(e^R)
 \ge c e^{(1-\sigma)R}
\]
for large \(R\). If \(\sigma<0\), the lower bound \(ce^R\) already
suffices. These estimates contradict polynomial growth. They also
exclude finite mass.
\end{proof}

The Euler product gives a transform on its actual convergence domain.
For \(\operatorname{Re}s>1\),
\begin{equation}
 \zeta(s)=\sum_{n\ge1}n^{-s}
         =\prod_p(1-p^{-s})^{-1},\qquad
 -\frac{\zeta'(s)}{\zeta(s)}
       =\sum_{n\ge2}\Lambda(n)n^{-s}.
 \label{v4-arith:weil:eq:euler}
\end{equation}
Absolute convergence proves the product and permits logarithmic
differentiation on compact subsets of this half-plane. Thus
\(-\zeta'/\zeta\) is the Laplace transform of \(\mu_0\) there.
With Fourier convention \(\int e^{-itx}\,d\mu(x)\), the transform of
\(\mu_\sigma\), \(\sigma>1\), is
\(-\zeta'(\sigma+it)/\zeta(\sigma+it)\).

The seminorm estimate in the theorem and dominated convergence give
\(\mu_{1+\varepsilon}\to\mu_1\) weakly in tempered distributions.
Fourier continuity gives the corresponding boundary transform.
Meromorphic continuation of \(-\zeta'/\zeta\) does not extend the original
Laplace integral beyond its convergence domain.

\section{Theta continuation and the zero divisor}
\label{v4-arith:weil:theta}

We normalize the completed zeta function before forming its logarithmic
derivative. Let
\[
 \theta(t)=\sum_{n\in\mathbb Z}e^{-\pi n^2t},\qquad
 r(t)=\frac{\theta(t)-1}{2}\qquad(t>0).
\]
The Gaussian transform gives the theta identity
\begin{equation}
 \theta(t)=t^{-1/2}\theta(t^{-1}).
 \label{v4-arith:weil:eq:theta}
\end{equation}
One proof periodizes \(e^{-\pi t x^2}\). Its \(m\)-th Fourier coefficient
is \(t^{-1/2}e^{-\pi m^2/t}\). Both periodic sums converge uniformly,
and uniqueness of Fourier coefficients gives the identity at \(x=0\).

The Gaussian transform itself follows by differentiating its integral
with respect to frequency and using its value at zero.

For \(\operatorname{Re}s>1\), termwise integration yields
\[
 \Lambda_\zeta(s):=\pi^{-s/2}\Gamma(s/2)\zeta(s)
             =\int_0^\infty r(t)t^{s/2}\,\frac{dt}{t}.
\]
Split the integral at one and use \eqref{v4-arith:weil:eq:theta}
in the lower part. The result is
\begin{equation}
 \Lambda_\zeta(s)=\frac1{s(s-1)}
    +\int_1^\infty r(t)
            \bigl(t^{s/2}+t^{(1-s)/2}\bigr)\,\frac{dt}{t}.
 \label{v4-arith:weil:eq:completed}
\end{equation}
The integral is entire because \(r(t)\) decreases exponentially.
It follows that \(\Lambda_\zeta\) is meromorphic, with residues
\(-1\) at zero and \(1\) at one.

\begin{definition}[The normalized entire function]
\label{v4-arith:weil:def:xi}
Set
\begin{equation}
 \xi(s)=\frac12s(s-1)\Lambda_\zeta(s)
 =\frac12+\frac{s(s-1)}2
    \int_1^\infty r(t)
          \bigl(t^{s/2}+t^{(1-s)/2}\bigr)\,\frac{dt}{t}.
 \label{v4-arith:weil:eq:xi}
\end{equation}
Its zeros are counted with their analytic multiplicities.
For a zero \(\rho\), put \(z_\rho=-i(\rho-\tfrac12)\).
\end{definition}

The integral proves
\[
 \xi(0)=\xi(1)=\frac12,\qquad
 \xi(1-s)=\xi(s),\qquad
 \xi(\bar s)=\overline{\xi(s)}.
\]
The Euler product has no zero for \(\operatorname{Re}s>1\).
The functional equation therefore places all zeros of \(\xi\) in
\(0\le\operatorname{Re}s\le1\). The associated parameters satisfy
\(|\operatorname{Im}z_\rho|\le1/2\).

There is no real zero of \(\xi\). For \(0<s<1\), pair consecutive
terms of the alternating series:
\[
 \eta(s)=\sum_{m\ge1}
       \bigl((2m-1)^{-s}-(2m)^{-s}\bigr)>0.
\]
Dirichlet convergence and analytic continuation give
\(\eta(s)=(1-2^{1-s})\zeta(s)\). Hence \(\zeta(s)\ne0\) on that
interval. The endpoints and the remaining real line follow from
the displayed values, the Euler product and reflection.
Consequently no \(z_\rho\) lies on the imaginary axis.

\begin{lemma}[Growth and summability]
\label{v4-arith:weil:lem:zeros}
For \(R\ge2\),
\[
 \log\max_{|s|\le R}|\xi(s)|\le CR\log(R+2).
\]
The number of zeros in \(|s|\le R\), with multiplicity, is
\(O(R\log(R+2))\). In particular,
\(\sum_\rho|\rho|^{-2}<\infty\), and the corresponding inverse-square
sum over \(z_\rho\) converges. There are infinitely many zeros.
\end{lemma}
\begin{proof}
On \(t\ge1\), one has \(r(t)\le C e^{-\pi t}\).
For \(|s|\le R\), the two powers in \eqref{v4-arith:weil:eq:xi}
are bounded by \(2t^{(R+1)/2}\).
The elementary maximum of \(t^A e^{-\pi t/2}\) is
\((2A/(\pi e))^A\). This bounds the integral by
\(\exp(CR\log(R+2))\), including its quadratic prefactor.

Jensen's formula at radius \(2R\), with \(\xi(0)=1/2\), bounds the
number of zeros in radius \(R\): each contributes at least \(\log2\).
Summing over dyadic annuli proves inverse-square summability.
Translation by \(1/2\) gives the assertion for \(z_\rho\).

If the zero set were finite, Hadamard factorization for an entire
function of order at most one would give \(\xi(s)=e^{as+b}P(s)\)
for a polynomial \(P\). Along positive real \(s\), Stirling's formula
and \(\zeta(s)\to1\) instead give
\(\log|\xi(s)|=\tfrac12s\log s+O(s)\).
This contradicts that representation.
\end{proof}

\begin{proposition}[The genus-one normalization]
\label{v4-arith:weil:prop:product}
Let \(\gamma_E\) be Euler's constant and put
\[
 B=-1-\frac{\gamma_E}{2}+\frac12\log(4\pi).
\]
Then
\begin{equation}
 \xi(s)=\frac12e^{Bs}
       \prod_\rho(1-s/\rho)e^{s/\rho},\qquad
 \frac{\xi'(s)}{\xi(s)}
       =B+\sum_\rho\left(\frac1{s-\rho}+\frac1\rho\right).
 \label{v4-arith:weil:eq:product}
\end{equation}
The factors and paired summands converge locally uniformly away from
their zeros or poles.
\end{proposition}
\begin{proof}
The preceding growth estimate permits Hadamard factorization of genus
one. Inverse-square summability gives local normal convergence.
At zero the constant factor is \(1/2\), and the logarithmic derivative
of every elementary factor vanishes. The linear exponential coefficient
is therefore \(\xi'(0)/\xi(0)\).

Use \(\Gamma(1+s/2)=(s/2)\Gamma(s/2)\) to write
\[
 \xi(s)=(s-1)\Gamma(1+s/2)\pi^{-s/2}\zeta(s).
\]
The values \(\zeta(0)=-1/2\), \(\zeta'(0)=-\tfrac12\log(2\pi)\)
and \(\Gamma'(1)=-\gamma_E\) give the stated \(B\).
For the two zeta values, expand the equivalent functional equation
\[
 \zeta(s)=2^s\pi^{s-1}\sin(\pi s/2)\Gamma(1-s)\zeta(1-s).
\]
The product \(\Gamma(1-s)\zeta(1-s)\) is \(-s^{-1}+O(s)\).
The remaining factor is
\(\tfrac s2(1+s\log(2\pi)+O(s^2))\).
This proves both values.

Local differentiation of the normally
convergent canonical product proves the logarithmic derivative formula.
\end{proof}

The summands in \eqref{v4-arith:weil:eq:product} are paired.
The formula makes no separate absolute-convergence assertion for
\(\sum_\rho1/\rho\). This convention will also determine the normalization
of spectral determinants.

\subsection{The argument term in zero counting}

Let \(N(T)\) count the zeros \(\rho\) with
\(0<\operatorname{Im}\rho\le T\), including multiplicity.
For \(T>0\) that is not a zero ordinate, define
\[
 \vartheta(T)=
 \operatorname{Im}\log\Gamma(\tfrac14+\tfrac{iT}{2})
                -\frac T2\log\pi,\qquad
 S(T)=\frac1\pi\Delta_{\,2\to2+iT\to1/2+iT}\arg\zeta.
\]
The gamma logarithm is the analytic branch on the right half-plane
that is real on the positive real axis. The zeta argument starts
at zero at \(s=2\).

\begin{proposition}[The zero-count normalization]
\label{v4-arith:weil:prop:zero-count}
Away from zero ordinates,
\begin{equation}
 N(T)=1+\frac{\vartheta(T)}\pi+S(T)
 =\frac{T}{2\pi}\log\frac{T}{2\pi e}
       +\frac78+S(T)+O(T^{-1}).
 \label{v4-arith:weil:eq:zero-count}
\end{equation}
At a zero ordinate, the formula for the stated convention \(N(T)\)
uses the right limit of \(S\). The classical bound
\(S(T)=O(\log T)\) gives
\[
 N(T)=\frac{T}{2\pi}\log\frac{T}{2\pi e}+O(\log T).
\]
\end{proposition}
\begin{proof}
Use the rectangle with real parts \(-1,2\) and imaginary parts
\(0,T\). There are no zeros on its boundary. The real bottom edge
contributes no argument change, since \(\xi\) is real and nonzero there.
Reflection \(s\mapsto1-\bar s\) and conjugation of \(\xi\) show that
the remaining argument change is twice the change along
\[
 C_T:\quad 2\longrightarrow2+iT\longrightarrow\tfrac12+iT.
\]
The argument principle therefore gives
\(N(T)=\pi^{-1}\Delta_{C_T}\arg\xi\).

At the endpoint, the continuous arguments of \(s\) and \(s-1\)
sum to \(\pi\). The factorization
\(\xi(s)=\tfrac12s(s-1)\pi^{-s/2}\Gamma(s/2)\zeta(s)\)
now gives the first equality. In particular, the horizontal
segment contributes to \(S(T)\).

Stirling's expansion at \(z=1/4+iT/2\) gives
\[
 \vartheta(T)=\frac T2\log\frac{T}{2\pi}
       -\frac T2-\frac\pi8+\frac1{48T}+O(T^{-3}).
\]
For the constant, the term
\((z-\tfrac12)\log z-z\) contributes
\(-\tfrac14\arg z\to-\pi/8\).
Thus \(7/8=1-1/8\).
Taking right limits proves the convention at zero heights.

The argument bound is the classical theorem. A precise primary
estimate is Hasanalizade--Shen--Wong,
\emph{Counting zeros of the Riemann zeta function},
Theorem~1.3 and Corollary~1.4,
\href{https://arxiv.org/abs/2107.06506}{arXiv:2107.06506}.
Their Corollary~1.2 also supplies the displayed estimate for \(N(T)\).
\end{proof}

This count controls ordinary convergent sums. It does not regularize
divergent positive-power moments.
 For example, put
\[
 M_n(T)=\sum_{0<\operatorname{Im}\rho\le T}
                 \bigl(\tfrac14+(\operatorname{Im}\rho)^2\bigr)^n.
\]
For every integer \(n\ge0\), partial summation gives
\begin{equation}
 M_n(T)=
 \frac{T^{2n+1}}{2\pi(2n+1)}\log\frac{T}{2\pi}
 -\frac{T^{2n+1}}{2\pi(2n+1)^2}
 +O_n(T^{2n}\log T).
 \label{v4-arith:weil:eq:positive-moments}
\end{equation}

Indeed, integrate \((1/4+t^2)^n\) against \(dN(t)\).
The error \(O(\log(2+t))\) in the count contributes
\(O_n(T^{2n}\log T)\). In the main integral, replacing
\((1/4+t^2)^n\) by \(t^{2n}\) gives a smaller error.
The integral of \(t^{2n}\log(t/(2\pi))/(2\pi)\) is the displayed main term.

Every \(M_n(T)\) diverges. A finite arithmetic moment functional
therefore requires an additional, compatible regularization.

\section{The complete arithmetic functional}
\label{v4-arith:weil:functional}

Write \(\psi_\Gamma=\Gamma'/\Gamma\) and define the real even functions
\begin{equation}
 a(t)=\operatorname{Re}\psi_\Gamma(\tfrac14+\tfrac{it}{2})-\log\pi,
 \qquad k(t)=-a(t).
 \label{v4-arith:weil:eq:gamma-kernels}
\end{equation}
They are smooth on the real line and have at most logarithmic growth.

For \(f\in\mathscr A_c\), set
\begin{align}
 P(f)&=h_f(i/2)+h_f(-i/2),\notag\\
 W_\infty(f)&=\frac1{2\pi}\int_{\mathbb R}h_f(t)k(t)\,dt,\notag\\
 W_{\mathrm{fin}}(f)&=
       \sum_{n\ge2}\frac{\Lambda(n)}{\sqrt n}
                    \bigl(f(\log n)+f(-\log n)\bigr),\notag\\
 W(f)&=P(f)-W_\infty(f)-W_{\mathrm{fin}}(f).
 \label{v4-arith:weil:eq:functional}
\end{align}
The prime sum is finite. The gamma integral is absolutely convergent.
These formulas define continuous complex-linear functionals on the
compact test space. The sign of \(W_\infty\) is fixed by its density \(k\).

\begin{theorem}[The complete explicit formula]
\label{v4-arith:weil:thm:explicit}
For every \(f\in\mathscr A_c\),
\begin{equation}
 W(f)=\sum_\rho h_f(z_\rho).
 \label{v4-arith:weil:eq:explicit}
\end{equation}
The sum over zeros is absolutely convergent, with multiplicity.
For \(f,g\in\mathscr A_c\), the form
\begin{equation}
 q_W(f,g)=W(f*g^*)
       =\sum_\rho h_f(z_\rho)\overline{h_g(\bar z_\rho)}
 \label{v4-arith:weil:eq:pairing}
\end{equation}
is Hermitian and linear in its first argument.
\end{theorem}
\begin{proof}
Put \(H(s)=\Phi_f(s)\), and choose \(1<c<3/2\).
Let
\[
 G(s)=\frac{\Lambda_\zeta'(s)}{\Lambda_\zeta(s)}
      =-\frac12\log\pi+\frac12\psi_\Gamma(s/2)
                         +\frac{\zeta'(s)}{\zeta(s)}.
\]
Inside the strip \(1-c\le\operatorname{Re}s\le c\), its residues
are the zero multiplicities and \(-1\) at each of \(0,1\).
The functional equation gives \(G(1-s)=-G(s)\).

We justify the contour limit. For all sufficiently large integers \(j\),
choose \(T_j\in[j,j+1]\) at distance at least \(j^{-2}\) from every zero
ordinate. The excluded intervals have total length
\(O(\log j/j)<1\), by Lemma~\ref{v4-arith:weil:lem:zeros}.

On the horizontal segments at heights \(\pm T_j\),
\eqref{v4-arith:weil:eq:product} gives a polynomial bound for
\(\xi'/\xi\). The zeros with \(|\rho|\le2T_j\) are bounded using the
chosen distance. For the others,
\[
 \left|\frac1{s-\rho}+\frac1\rho\right|
       \le C\frac{|s|}{|\rho|^2}.
\]
The rational correction
\(G=\xi'/\xi-1/s-1/(s-1)\) has the same bound.
Rapid vertical-strip decay of \(H\) makes both horizontal integrals
tend to zero.

Apply the residue theorem to \(G(s)H(s)\) on that rectangle.
Reflection of the left vertical side gives
\begin{equation}
 \sum_\rho H(\rho)-H(0)-H(1)
 =\frac1{2\pi i}\int_{\operatorname{Re}s=c}
             G(s)\bigl(H(s)+H(1-s)\bigr)\,ds.
 \label{v4-arith:weil:eq:contour}
\end{equation}
Absolute convergence of the zero sum follows from strip decay and
the zero count.

On the right line, the Dirichlet series in
\eqref{v4-arith:weil:eq:euler} is absolutely convergent.
Termwise integration is permitted. Fourier inversion gives
\[
 \frac1{2\pi i}\int_{\operatorname{Re}s=c}
                 n^{-s}H(s)\,ds=n^{-1/2}f(\log n).
\]
The reflected test gives \(n^{-1/2}f(-\log n)\).
The negative sign in \(\zeta'/\zeta\) therefore produces
\(-W_{\mathrm{fin}}(f)\).

Move the gamma part of the integral to \(\operatorname{Re}s=1/2\).
Its logarithmic derivative has no pole between these lines.
The two reflected integrals combine to
\[
 \frac1{2\pi}\int_{\mathbb R}h_f(t)
       \bigl(\operatorname{Re}\psi_\Gamma(\tfrac14+\tfrac{it}{2})
                     -\log\pi\bigr)\,dt=-W_\infty(f).
\]
Finally, \(H(0)=h_f(i/2)\) and \(H(1)=h_f(-i/2)\).
This proves \eqref{v4-arith:weil:eq:explicit}.

Fourier multiplication proves the paired formula.
The zero multiset is invariant under
\(\rho\mapsto1-\bar\rho\), which sends \(z_\rho\) to \(\bar z_\rho\).
Reindexing the absolutely convergent sum proves Hermitian symmetry.
\end{proof}

In \eqref{v4-arith:weil:eq:pairing}, conjugation acts at
\(\bar z_\rho\). Replacing the summand by \(|h_f(z_\rho)|^2\)
would change the form when \(z_\rho\) is nonreal.
Under the Riemann hypothesis every \(z_\rho\) is real, and that
specialization is valid.

\section{Exponential extensions and their topology}
\label{v4-arith:weil:exponential}

Gaussian tests require a domain larger than \(\mathscr A_c\).
The growth threshold in Theorem~\ref{v4-arith:weil:thm:prime-threshold}
indicates the needed exponential margin.

\begin{definition}[Exponential test spaces]
\label{v4-arith:weil:def:exponential}
For \(a_0>1/2\), let
\begin{equation}
 \mathscr A_{a_0}=
 \left\{f\in C^\infty(\mathbb R):
 p_{a_0,j}(f):=\int_{\mathbb R}
          e^{a_0|x|}|f^{(j)}(x)|\,dx<\infty
          \text{ for every }j\ge0\right\}.
 \label{v4-arith:weil:eq:exponential}
\end{equation}
Give this space the topology of its displayed seminorms.
Put
\[
 \mathscr A_{\exp}=\bigcup_{a_0>1/2}\mathscr A_{a_0},\qquad
 \mathscr A_\infty=\bigcap_{a_0>1/2}\mathscr A_{a_0}.
\]
The union has the final locally convex topology of the steps
\(\mathscr A_{1/2+1/n}\). The intersection has all the exponential
seminorms. Convolution and the involution are as in
\eqref{v4-arith:weil:eq:convolution}.
\end{definition}

\begin{theorem}[Continuity on each exponential step]
\label{v4-arith:weil:thm:extension}
Each \(\mathscr A_{a_0}\) is a complete Fr\'echet convolution
\(*\)-algebra, and \(\mathscr A_c\) is dense in it.
The transform \(h_f\) is holomorphic for
\(|\operatorname{Im}z|<a_0\). All its derivatives decrease rapidly
on smaller closed strips.

Every term of \eqref{v4-arith:weil:eq:functional} extends continuously
to this step. There is a constant \(C_{a_0}\) such that
\begin{equation}
 |W(f)|\le C_{a_0}
       \bigl(p_{a_0,0}(f)+p_{a_0,1}(f)+p_{a_0,2}(f)\bigr).
 \label{v4-arith:weil:eq:continuity}
\end{equation}
The explicit formula and its Hermitian pairing remain valid, with
absolutely convergent prime and zero sums.
\end{theorem}
\begin{proof}
The inequality \(e^{a_0|x+y|}\le e^{a_0|x|}e^{a_0|y|}\) gives
\[
 p_{a_0,j}(f*g)\le p_{a_0,j}(f)p_{a_0,0}(g).
\]
The involution preserves each seminorm. A Cauchy sequence has weighted
\(L^1\) limits for every derivative. Distributional differentiation
identifies successive limits. Local one-dimensional Sobolev regularity
then gives a smooth representative with those derivatives.
This proves completeness.

Choose \(\chi\in C_c^\infty(\mathbb R)\) equal to one near zero.
Leibniz's rule shows
\(\chi(x/R)f(x)\to f(x)\) in every seminorm.
The errors are integrable derivative tails multiplied by bounded
cutoff derivatives and nonpositive powers of \(R\).
This proves density.

On a strip \(|\operatorname{Im}z|\le b<a_0\), the spare weight
\(e^{(a_0-b)|x|}\) controls every polynomial in \(x\).
Differentiation in \(z\) and integration by parts in \(x\) give the
asserted holomorphy and rapid decay. The two polar evaluations are
bounded by \(p_{a_0,0}(f)\).
On the real line,
\[
 |h_f(t)|\le
 \min\{p_{a_0,0}(f),\,|t|^{-2}p_{a_0,2}(f)\}.
\]
Logarithmic growth of \(k\) therefore controls the gamma integral.

The function \(e^{a_0|x|}f(x)\) has an integrable weak derivative.
Integrating from infinity on either half-line gives
\[
 \sup_x e^{a_0|x|}|f(x)|
       \le a_0p_{a_0,0}(f)+p_{a_0,1}(f).
\]
Consequently the absolute prime sum is at most
\[
 2\bigl(a_0p_{a_0,0}(f)+p_{a_0,1}(f)\bigr)
           \sum_{n\ge2}\frac{\log n}{n^{a_0+1/2}}.
\]
This series converges because \(a_0>1/2\), proving
\eqref{v4-arith:weil:eq:continuity}.

Strip decay and the zero count similarly control the zero sum.
Pass to the limit through the compact cutoffs to obtain the explicit
formula. Convolution continuity proves the paired assertion.
\end{proof}

The extensions agree on overlapping steps, so they define \(W\) on
\(\mathscr A_{\exp}\). All assertions about passage of a quadratic
inequality below use a specified step or the compact test topology.
No completeness claim for the union is needed.

The intersection \(\mathscr A_\infty\) is the space of smooth functions
with rapid exponential decay. Its topology is equivalently given by
\begin{equation}
 \sup_x e^{m|x|}(1+|x|)^N|f^{(j)}(x)|,
       \qquad m,N,j\in\mathbb N.
 \label{v4-arith:weil:eq:rapid-exponential}
\end{equation}
Indeed, one extra exponential weight bounds a polynomial weight and
an integral. Conversely, the weighted \(L^1\) bounds for two consecutive
derivatives bound the weighted supremum as in the proof above.
Compact cutoffs converge in all these seminorms.

The gamma functional alone extends continuously to Schwartz space.
The prime functional does not have such an extension: its positive
measure contains \(\mu_{1/2}\) and its reflection.
Thus Schwartz approximation alone does not justify passing each
arithmetic term to a limit.

\begin{proposition}[The regular convolution norm]
\label{v4-arith:weil:prop:regular}
Convolution by \(f\in\mathscr A_{\exp}\) acts boundedly on
\(L^2(\mathbb R)\), with operator norm \(\|h_f\|_\infty\).
The completion for this norm is \(C_0(\mathbb R)\) under Fourier
transform.
\end{proposition}
\begin{proof}
Plancherel turns convolution into multiplication by \(h_f\).
Since \(f\in L^1\), its transform is continuous and vanishes at infinity.
To prove density, start with \(H\in C_c^\infty(\mathbb R)\) in the
frequency variable. Its inverse Fourier transform \(v\) is Schwartz.
Compact cutoffs \(v_R\) converge to \(v\) in \(L^1\), so
\(\|h_{v_R}-H\|_\infty\le\|v_R-v\|_1\to0\).
Finally, \(C_c^\infty\) is uniformly dense in \(C_0\).
\end{proof}

\section{The polar and finite-prime forms}
\label{v4-arith:weil:signed-pieces}

For \(f\in\mathscr A_{\exp}\), write
\[
 \ell_-(f)=\int f(x)e^{-x/2}\,dx,\qquad
 \ell_+(f)=\int f(x)e^{x/2}\,dx.
\]
Both are continuous on every exponential step.

\begin{proposition}[The full polar matrix]
\label{v4-arith:weil:prop:polar}
One has
\begin{equation}
 P(f*g^*)=
 \ell_-(f)\overline{\ell_+(g)}
       +\ell_+(f)\overline{\ell_-(g)}.
 \label{v4-arith:weil:eq:polar}
\end{equation}
On compact tests, the map
\((\ell_-,\ell_+):\mathscr A_c\to\mathbb C^2\) is surjective.
The polar form has rank two, with one positive and one negative
direction. Its kernel is the common kernel of the two moments.
On the complex even tests it is the positive form
\(2\ell_+(f)\overline{\ell_+(g)}\).
\end{proposition}
\begin{proof}
Evaluate \(h_fh_g^\sharp\) at \(i/2\) and \(-i/2\).
A linear relation between the two moment functionals would make
\(c_-e^{-x/2}+c_+e^{x/2}\) vanish identically.
Its value and derivative at zero force \(c_-=c_+=0\).
The map therefore has two-dimensional range.

The matrix on that range is
\(\left(\begin{smallmatrix}0&1\\1&0\end{smallmatrix}\right)\).
Its eigenvalues are \(1,-1\). Evenness gives \(\ell_-=\ell_+\),
which proves the last assertion.
\end{proof}

A right inverse \(S:\mathbb C^2\to\mathscr A_c\) gives the continuous
projection \(1-S(\ell_-,\ell_+)\) onto the polar kernel.
A single vector with norm square \(|\ell_+(f)|^2\) does not realize
the full polar form.

The prime term has a different finite-support realization.
For \(R>0\), let \(j_R:L^2([-R,R])\to L^2(\mathbb R)\) extend by zero.
Write \(T_a v(x)=v(x-a)\), and define
\begin{equation}
 B_R=\sum_{2\le n\le e^{2R}}\frac{\Lambda(n)}{\sqrt n}
                  j_R^*(T_{\log n}+T_{-\log n})j_R.
 \label{v4-arith:weil:eq:prime-operator}
\end{equation}

\begin{proposition}[The finite-prime operator]
\label{v4-arith:weil:prop:prime-operator}
The operator \(B_R\) is bounded and self-adjoint, with
\(\|B_R\|\le2\sum_{n\le e^{2R}}\Lambda(n)n^{-1/2}\).
For smooth \(f,g\) supported in \((-R,R)\),
\[
 W_{\mathrm{fin}}(f*g^*)=\langle B_Rf,g\rangle
 =\langle |B_R|^{1/2}f,
                  \operatorname{sgn}(B_R)|B_R|^{1/2}g\rangle.
\]
These pairings agree under enlargement of the support interval.
\end{proposition}
\begin{proof}
Translations are unitary and \(T_a^*=T_{-a}\), proving the operator
claims and the norm bound. A change of variables gives
\[
 (f*g^*)(a)=\int f(x)\overline{g(x-a)}\,dx
                     =\langle T_{-a}j_Rf,j_Rg\rangle.
\]
This vanishes for \(|a|\ge2R\). Summing at both signs of each
prime-power logarithm proves the first identity and compatibility.
The spectral theorem for \(B_R\), with \(\operatorname{sgn}(0)=0\),
proves the second identity.
\end{proof}

\begin{example}[Two translated bumps]
\label{v4-arith:weil:ex:prime-signs}
Put \(L=\log2\). Choose a nonzero real smooth \(b\) supported in
\((-\varepsilon,\varepsilon)\), where
\(2\varepsilon<\min\{\log2,\log(3/2)\}\).
On the span of \(b,T_Lb\), the finite-prime form has matrix
\begin{equation}
 \frac{\log2}{\sqrt2}\|b\|_2^2
                  \begin{pmatrix}0&1\\1&0\end{pmatrix}.
 \label{v4-arith:weil:eq:prime-signs}
\end{equation}
Indeed, \(b*b^*\) is supported in \([-2\varepsilon,2\varepsilon]\)
and has value \(\|b\|_2^2\) at zero. Only the two shifts
\(\pm\log2\) meet the off-diagonal correlations.

Hence
\[
 W_{\mathrm{fin}}\bigl((b\pm T_Lb)*(b\pm T_Lb)^*\bigr)
                    =\pm\sqrt2\log2\,\|b\|_2^2.
\]
This excludes a positive Hilbert realization of the finite-prime
form on all compact tests. In the full form that term is subtracted.
Its negative value is therefore not a negative full Weil square.
\end{example}

Prime powers remain essential even in such small examples.
Two sufficiently narrow bumps separated by \(\log4\) have an
off-diagonal contribution \(\Lambda(4)/\sqrt4=(\log2)/2\).
Retaining only primes would erase that contribution.

\section{The gamma kernel and its auxiliary form}
\label{v4-arith:weil:gamma}

The signed gamma density \(k\) has both signs. Its absolute value
cannot be absorbed into an unspecified positive measure without
changing the arithmetic functional.
To see its shape explicitly, set \(c_n=n+1/4\). The digamma series gives
\begin{equation}
 a(t)=a(0)+\sum_{n\ge0}
       \frac{t^2/4}{c_n(c_n^2+t^2/4)},\qquad
 a(0)=-\gamma_E-\frac\pi2-3\log2-\log\pi.
 \label{v4-arith:weil:eq:digamma}
\end{equation}
Here and below \(a=-k\), as in
\eqref{v4-arith:weil:eq:gamma-kernels}.
See
\href{https://dlmf.nist.gov/5.7.E6}{DLMF 5.7.6} and
\href{https://dlmf.nist.gov/5.4.E19}{5.4.19}.

\begin{proposition}[The gamma signs]
\label{v4-arith:weil:prop:gamma-signs}
The function \(a\) is even and strictly increasing on \((0,\infty)\).
It has its minimum at zero and satisfies
\[
 a(t)=\log\frac{|t|}{2\pi}+O(t^{-2})
                  \qquad(|t|\to\infty).
\]
Thus \(k(0)>0\), and \(k(t)\to-\infty\).
The function \(|\Gamma(1/4+it/2)|^2\) is not a constant multiple of
\(k(t)\).
\end{proposition}
\begin{proof}
Subtract the digamma series at zero. This gives
\eqref{v4-arith:weil:eq:digamma}, whose summands are nonnegative and
strictly increasing in \(|t|\). Their derivatives converge uniformly
on compact intervals:
\[
 a'(t)=\sum_{n\ge0}
       \frac{c_nt/2}{(c_n^2+t^2/4)^2}>0\qquad(t>0).
\]
The value at zero follows from the reflection and duplication identities
for \(\psi_\Gamma\). Stirling's expansion
\(\psi_\Gamma(z)=\log z-1/(2z)+O(|z|^{-2})\), on the vertical line
\(\operatorname{Re}z=1/4\), gives the asymptotic.
Finally, the gamma absolute square is strictly positive everywhere,
whereas \(k\) changes sign.
\end{proof}

The phase and the absolute square also have different derivatives:
\begin{equation}
 \frac{d}{dt}\arg\Gamma(\tfrac14+\tfrac{it}{2})
       =\frac12\operatorname{Re}\psi_\Gamma(\tfrac14+\tfrac{it}{2}),
 \qquad
 \frac{d}{dt}\log|\Gamma(\tfrac14+\tfrac{it}{2})|^2
       =-\operatorname{Im}\psi_\Gamma(\tfrac14+\tfrac{it}{2}).
 \label{v4-arith:weil:eq:gamma-derivatives}
\end{equation}
These follow by taking the imaginary and real parts of
\((i/2)\psi_\Gamma\). A local continuous phase is sufficient.
For \(\Lambda_\infty(s)=\pi^{-s/2}\Gamma(s/2)\), the arithmetic
density is
\[
 k(t)=-2\operatorname{Re}
      \frac{\Lambda_\infty'}{\Lambda_\infty}(\tfrac12+it).
\]
This specifies the operation and normalization that a local
archimedean comparison must preserve.

\begin{proposition}[The kernel away from the origin]
\label{v4-arith:weil:prop:off-origin}
On compact tests supported away from zero, the distribution
\(P-W_\infty\) is integration against
\begin{equation}
 K(x)=2\cosh(x/2)-\frac{e^{-|x|/2}}{1-e^{-2|x|}},
                      \qquad x\ne0.
 \label{v4-arith:weil:eq:off-origin}
\end{equation}
In particular, if \(0<b<c<\log2\) and \(f\ne0\) is even,
nonnegative and supported in \((-c,-b)\cup(b,c)\), then
\[
 W_{\mathrm{fin}}(f)=0,\qquad
 W_\infty(f)=2\int_b^c
              \frac{e^{-x/2}}{1-e^{-2x}}f(x)\,dx>0.
\]
\end{proposition}
\begin{proof}
The polar functional is integration against \(2\cosh(x/2)\).
Use the integral identity
\[
 \psi_\Gamma(z)+\gamma_E
       =\int_0^\infty\frac{e^{-v}-e^{-zv}}{1-e^{-v}}\,dv,
                       \qquad\operatorname{Re}z>0,
\]
as in \href{https://dlmf.nist.gov/5.9.E16}{DLMF 5.9.16}.
For a test vanishing near zero, constant Fourier terms contribute
zero. Fourier inversion of the cosine term, followed by \(v=2|x|\),
gives the second term of \eqref{v4-arith:weil:eq:off-origin}
with its negative sign. Truncating the \(v\)-integral first justifies
the calculation.

The support condition removes its singular endpoint,
and exponential decay controls infinity.
The final assertion follows because no prime-power logarithm lies
in the stated support.
\end{proof}

\subsection{A second gamma density}

Define
\[
 b(t)=\operatorname{Re}\psi_\Gamma(\tfrac14+\tfrac{it}{2})+\log\pi
                   =2\log\pi-k(t),
\]
and the auxiliary form
\begin{equation}
 G_+(f,g)=\frac1{2\pi}\int_{\mathbb R}
                  h_f(t)\overline{h_g(t)}b(t)\,dt.
 \label{v4-arith:weil:eq:auxiliary}
\end{equation}
It is continuous on Schwartz space. Plancherel gives
\begin{equation}
 W_\infty(f*g^*)=
          2\log\pi\,\langle f,g\rangle_{L^2}-G_+(f,g).
 \label{v4-arith:weil:eq:auxiliary-comparison}
\end{equation}
The function \(b\) is strictly increasing in \(|t|\), with
\[
 b(0)=-\gamma_E-\frac\pi2-3\log2+\log\pi<0,\qquad
 b(t)=\log(|t|\pi/2)+O(t^{-2}).
\]
It therefore has a unique positive zero, denoted \(t_b\).

The interval \(0.84<t_b<0.86\) follows from a positive-tail estimate.
Put
\[
 S_N(t)=b(0)+\sum_{n=0}^N
          \frac{t^2/4}{c_n(c_n^2+t^2/4)}.
\]
Then
\[
 0\le b(t)-S_N(t)\le\frac{t^2}{8(N+1/4)^2}.
\]
At \(N=20\), the rational bounds
\[
 \begin{gathered}
 0.577<\gamma_E<0.578,\qquad3.1415<\pi<3.1416,\\
 0.6931<\log2<0.6932,\qquad1.1447<\log\pi<1.1448
 \end{gathered}
\]
give
\[
 -0.024<S_{20}(0.84)<-0.020,\qquad
 0.017<S_{20}(0.86)<0.020.
\]
Both tails are less than \(0.00023\). The signs and monotonicity
prove the interval.

For a compact test, define the two nonnegative energies
\[
 A_b(f)=\int_{|t|<t_b}|h_f(t)|^2(-b(t))\,dt,\qquad
 B_b(f)=\int_{|t|>t_b}|h_f(t)|^2b(t)\,dt.
\]
Their difference gives \(G_+(f,f)=(B_b(f)-A_b(f))/(2\pi)\).
The sufficient high-frequency condition is
\begin{equation}
 2t_b|b(0)|\sup_{|t|\le t_b}|h_f(t)|^2\le B_b(f).
 \label{v4-arith:weil:eq:high-frequency}
\end{equation}
It bounds \(A_b(f)\) by \(B_b(f)\), and hence implies \(G_+(f,f)\ge0\).

\begin{proposition}[Compact high-frequency tests]
\label{v4-arith:weil:prop:high-frequency}
The class defined by \eqref{v4-arith:weil:eq:high-frequency} contains
nonzero real even compact tests. It is a proper subset of
\(\mathscr A_c\).
\end{proposition}
\begin{proof}
Choose a real even compact smooth \(\chi\) with \(\int\chi\ne0\).
For \(M>0\), let \(f_M(x)=\chi(x)\cos(Mx)\).
Then
\[
 h_{f_M}(t)=\tfrac12\bigl(h_\chi(t+M)+h_\chi(t-M)\bigr).
\]
Schwartz decay makes the left side of
\eqref{v4-arith:weil:eq:high-frequency} tend to zero faster than
every inverse power of \(M\).

Near \(t=M\), the translate \(h_\chi(t-M)\) has modulus bounded below
on a fixed interval, while \(h_\chi(t+M)\to0\) uniformly there.
Since \(b(t)\sim\log M\) on that interval, \(B_b(f_M)\ge c\log M\)
for large \(M\). This proves the first assertion.

For the second, set \(v_L(x)=L^{-1/2}\chi(x/L)\).
A change of variables gives
\[
 G_+(v_L,v_L)=\frac1{2\pi}\int_{\mathbb R}
              |h_\chi(t)|^2b(t/L)\,dt
         \longrightarrow b(0)\|\chi\|_2^2<0.
\]
Logarithmic growth and Schwartz decay justify dominated convergence.
Thus \(v_L\) fails \eqref{v4-arith:weil:eq:high-frequency}
for all sufficiently large \(L\).
\end{proof}

The relation to full Weil positivity retains the arithmetic terms.
Put \(R(f)=P(f*f^*)-W_{\mathrm{fin}}(f*f^*)\). Then
\begin{equation}
 q_W(f,f)=R(f)-2\log\pi\,\|f\|_2^2
                       +\frac{B_b(f)-A_b(f)}{2\pi}.
 \label{v4-arith:weil:eq:full-energy}
\end{equation}
Thus \(q_W(f,f)\ge0\) is equivalent to
\[
 \frac{A_b(f)}{2\pi}+2\log\pi\,\|f\|_2^2
       \le\frac{B_b(f)}{2\pi}+R(f).
\]
Auxiliary positivity controls only the difference of the two energies.

A nonnegative majorant \(v\ge-b\) on \([-t_b,t_b]\) gives the useful
bound \(A_b(f)\le\int_{-t_b}^{t_b}|h_f|^2v\).
For support in \([-R,R]\), Cauchy--Schwarz gives
\[
 \sup_t|h_f(t)|^2\le2R\|f\|_2^2,\qquad
 A_b(f)\le4t_bR|b(0)|\|f\|_2^2.
\]
These are sufficient estimates for the corresponding terms.
They place no bound on \(R(f)\).

\section{Gaussian normalization and cross terms}
\label{v4-arith:weil:gaussians}

All derivatives of a Gaussian have rapid exponential decay.
Theorem~\ref{v4-arith:weil:thm:extension} therefore applies to the
entire Gaussian family and to its convolution products.
For \(\lambda>0\), set
\[
 g_\lambda(x)=e^{-\lambda x^2/2}.
\]
Completing the square gives
\begin{equation}
 h_{g_\lambda}(t)=\sqrt{\frac{2\pi}{\lambda}}e^{-t^2/(2\lambda)},
 \qquad
 (g_\lambda*g_\mu)(x)=
 \sqrt{\frac{2\pi}{\lambda+\mu}}\,
         e^{-\lambda\mu x^2/[2(\lambda+\mu)]}.
 \label{v4-arith:weil:eq:gaussian-convolution}
\end{equation}
These identities include all normalization factors.

Define the normalized signed gamma integral
\begin{equation}
 H(\lambda)=\frac1{\sqrt{\pi\lambda}}
                   \int_{\mathbb R}e^{-t^2/\lambda}k(t)\,dt.
 \label{v4-arith:weil:eq:H}
\end{equation}
If \(\nu=2\lambda\mu/(\lambda+\mu)\), then
\begin{equation}
 W_\infty(g_\lambda*g_\mu)
       =\sqrt{\frac{2\pi}{\lambda+\mu}}\,H(\nu).
 \label{v4-arith:weil:eq:gaussian-gamma}
\end{equation}
This follows by inserting the Fourier product in the integral
with coefficient \(1/(2\pi)\).

\begin{theorem}[A convergent Gaussian series]
\label{v4-arith:weil:thm:erfc}
Let
\(\operatorname{erfc}(x)=\frac2{\sqrt\pi}\int_x^\infty e^{-u^2}\,du\).
For every \(\lambda>0\),
\begin{equation}
 H(\lambda)=k(0)-\sum_{n\ge0}
 \left(\frac1{c_n}
       -\frac{2\sqrt\pi}{\sqrt\lambda}
        e^{4c_n^2/\lambda}\operatorname{erfc}(2c_n/\sqrt\lambda)\right).
 \label{v4-arith:weil:eq:erfc}
\end{equation}
Each summand is nonnegative and at most \(\lambda/(8c_n^3)\).
The series converges uniformly on compact \(\lambda\)-intervals.
\end{theorem}
\begin{proof}
The identity \eqref{v4-arith:weil:eq:digamma} expresses
\(k(0)-k(t)\) as a sum of nonnegative functions.
Tonelli permits integration of this sum against the Gaussian.
For \(d>0\),
\[
 \int_{\mathbb R}\frac{e^{-t^2/\lambda}}{t^2+d^2}\,dt
 =\sqrt\pi\int_0^\infty
            \frac{e^{-d^2u}}{\sqrt{u+1/\lambda}}\,du
 =\frac\pi d e^{d^2/\lambda}\operatorname{erfc}(d/\sqrt\lambda).
\]
The first equality uses
\((t^2+d^2)^{-1}=\int_0^\infty e^{-u(t^2+d^2)}\,du\).
The substitution \(v=d\sqrt{u+1/\lambda}\) gives the second.
Set \(d=2c_n\) to obtain \eqref{v4-arith:weil:eq:erfc}.

Replacing \(c_n^2+t^2/4\) by \(c_n^2\) bounds its summand by
\[
 \frac1{\sqrt{\pi\lambda}}
       \int e^{-t^2/\lambda}\frac{t^2}{4c_n^3}\,dt
             =\frac{\lambda}{8c_n^3}.
\]
This proves uniform convergence.
\end{proof}

\begin{proposition}[The signed Gaussian threshold]
\label{v4-arith:weil:prop:gaussian-threshold}
The function \(H\) is continuous and strictly decreasing on
\((0,\infty)\). Its limits at zero and infinity are \(k(0)>0\)
and \(-\infty\), respectively. It has a unique zero \(\lambda_c>0\).
Consequently,
\[
 W_\infty(g_\lambda*g_\lambda)\ge0
                    \quad\Longleftrightarrow\quad\lambda\le\lambda_c.
\]
For \(a_j\ge0\) and \(0<\lambda_j\le\lambda_c\),
\[
 W_\infty\left(
       \left(\sum_j a_jg_{\lambda_j}\right)*
       \left(\sum_j a_jg_{\lambda_j}\right)^*\right)\ge0.
\]
\end{proposition}
\begin{proof}
After \(t=\sqrt\lambda\,u\), \(H(\lambda)\) is the average of
\(k(\sqrt\lambda\,u)\) against \(\pi^{-1/2}e^{-u^2}\,du\).
Strict radial decrease of \(k\) gives strict decrease of \(H\).
Logarithmic growth supplies an integrable majorant on compact
parameter intervals and near zero.
Monotone convergence applied to
\(k(0)-k(\sqrt\lambda\,u)\) gives divergence to infinity as
\(\lambda\to\infty\). This proves the threshold assertion.

If both \(\lambda,\mu\le\lambda_c\), their harmonic mean
\(2\lambda\mu/(\lambda+\mu)\) is at most \(\lambda_c\).
Every cross pairing in \eqref{v4-arith:weil:eq:gaussian-gamma}
is therefore nonnegative. Expand the finite sum with
coefficients \(a_i a_j\ge0\).

\end{proof}

The polar and prime values on a Gaussian square are
\begin{equation}
 \begin{aligned}
 P(g_\lambda*g_\lambda)&=\frac{4\pi}{\lambda}e^{1/(4\lambda)},\\
 W_{\mathrm{fin}}(g_\lambda*g_\lambda)
   &=2\sqrt{\frac\pi\lambda}
       \sum_{n\ge2}\frac{\Lambda(n)}{\sqrt n}
                          e^{-\lambda(\log n)^2/4}.
 \end{aligned}
 \label{v4-arith:weil:eq:gaussian-arithmetic}
\end{equation}
The sum converges for every positive \(\lambda\).
The full form subtracts both the prime and signed gamma terms
from the polar term. Positivity of the gamma cone alone gives
no sign for that complete expression.

\subsection{The heat normalization}

For \(t>0\), let
\[
 p_t(x)=\frac1{\sqrt{4\pi t}}e^{-x^2/(4t)},\qquad h_{p_t}(z)=e^{-tz^2}.
\]
The explicit formula becomes
\begin{equation}
 \begin{aligned}
 \sum_\rho e^{-tz_\rho^2}
 ={}&2e^{t/4}-\frac{\log\pi}{\sqrt{4\pi t}}
    +\frac1{2\pi}\int_{\mathbb R}
         e^{-tu^2}\operatorname{Re}\psi_\Gamma(\tfrac14+\tfrac{iu}{2})\,du\\
 &-\frac1{\sqrt{\pi t}}
       \sum_{n\ge2}\frac{\Lambda(n)}{\sqrt n}
                           e^{-(\log n)^2/(4t)}.
 \end{aligned}
 \label{v4-arith:weil:eq:heat}
\end{equation}
Each expression is ordinarily convergent.

On a compact interval of
positive \(t\), every parameter derivative has an integrable or summable
majorant. For the prime sum, a Gaussian in \(\log n\) dominates
every inverse power of \(n\). For the zero sum,
\(|e^{-tz_\rho^2}|\le e^{t/4}e^{-t(\operatorname{Im}\rho)^2}\).
The zero count proves its summability.

The difference between the two sides of \eqref{v4-arith:weil:eq:heat}
is zero by the explicit formula. Its vanishing is distinct from
nonnegativity of \(q_W\). The convolution law
\(p_s*p_t=p_{s+t}\) evaluates cross terms, but does not determine
their sign.

\subsection{Changing the Fourier scale}

For the convention
\(\widehat f_{2\pi}(u)=\int f(x)e^{-2\pi iux}\,dx\), define
\[
 (\mathcal R f)(x)=\frac1{2\pi}f\left(-\frac{x}{2\pi}\right).
\]
A change of variables gives
\begin{equation}
 h_{\mathcal Rf}=\widehat f_{2\pi},\qquad
 \mathcal R(f*g)=\mathcal Rf*\mathcal Rg,\qquad
 \mathcal R(f^*)=(\mathcal Rf)^*.
 \label{v4-arith:weil:eq:rescaling}
\end{equation}
If \(f\) has integrable derivatives weighted by \(e^{b_0|x|}\)
with \(b_0>\pi\), then \(\mathcal Rf\in\mathscr A_{b_0/(2\pi)}\).
Thus the complete arithmetic functional in this convention is
\[
 W_{2\pi}(f)=W(\mathcal Rf).
\]
Its prime term is
\[
 -\frac1{2\pi}\sum_{n\ge2}\frac{\Lambda(n)}{\sqrt n}
    \left(f\left(\frac{\log n}{2\pi}\right)
             +f\left(-\frac{\log n}{2\pi}\right)\right),
\]
and its gamma integral is
\((2\pi)^{-1}\int\widehat f_{2\pi}(u)a(u)\,du\).
Both the sampling points and the measure factor change.

For \(f_\ell(x)=e^{-\pi x^2/\ell}\), one has
\(\widehat f_{\ell,2\pi}(u)=\sqrt\ell e^{-\pi\ell u^2}\).
Multiplying \eqref{v4-arith:weil:eq:heat} at \(t=\pi\ell\)
by \(\sqrt\ell\) gives
\[
 \begin{aligned}
 \sqrt\ell\sum_\rho e^{-\pi\ell z_\rho^2}
 ={}&2\sqrt\ell e^{\pi\ell/4}-\frac{\log\pi}{2\pi}
 +\frac{\sqrt\ell}{2\pi}\int e^{-\pi\ell u^2}
                  \operatorname{Re}\psi_\Gamma(\tfrac14+\tfrac{iu}{2})\,du\\
 &-\frac1\pi\sum_{n\ge2}\frac{\Lambda(n)}{\sqrt n}
                         e^{-(\log n)^2/(4\pi\ell)}.
 \end{aligned}
\]
This also fixes the constants in the two-scale Gaussian identity
\(f_\ell*f_m=\sqrt{\ell m/(\ell+m)}\,f_{\ell+m}\).

\subsection{The auxiliary Gaussian cone}

The auxiliary density \(b\) has its own Gaussian average. Put
\[
 I(L)=\int_{\mathbb R}e^{-2\pi u^2/L}b(u)\,du,\qquad m_L=\sqrt{L/2}.
\]
Then
\begin{equation}
 \frac{I(L)}{m_L}=2\log\pi-H\left(\frac{L}{2\pi}\right).
 \label{v4-arith:weil:eq:auxiliary-H}
\end{equation}
The normalized integral is strictly increasing and has a unique
zero \(L_b>0\). Its equation is
\(H(L_b/(2\pi))=2\log\pi\).
This is a different level from \(H(\lambda_c)=0\).

The individual integrals
\[
 J_n(L)=\int e^{-2\pi u^2/L}\frac{c_n}{c_n^2+u^2/4}\,du
       =2\pi e^{8\pi c_n^2/L}
                 \operatorname{erfc}\bigl(2c_n\sqrt{2\pi/L}\bigr)
\]
give the absolutely convergent paired series
\[
 I(L)=m_L(-\gamma_E+\log\pi)
               +\sum_{n\ge0}\left(\frac{m_L}{n+1}-J_n(L)\right).
\]
Subtracting the value of the digamma series at zero improves its tail:
\begin{equation}
 \frac{I(L)}{m_L}=b(0)+
      \sum_{n\ge0}\left(\frac1{c_n}-\frac{J_n(L)}{m_L}\right),
 \qquad
 0\le\frac1{c_n}-\frac{J_n(L)}{m_L}
                     \le\frac{L}{16\pi c_n^3}.
 \label{v4-arith:weil:eq:auxiliary-series}
\end{equation}
The tail starting at \(n=N\ge1\) is at most
\(L/[32\pi(N-3/4)^2]\).
These formulas follow from Theorem~\ref{v4-arith:weil:thm:erfc}
with \(\lambda=L/(2\pi)\). The difference
\(1/(n+1)-1/c_n\) is absolutely summable, which justifies regrouping.

The same real Gaussian integral also gives, for \(\alpha,d>0\),
\[
 \int_{\mathbb R}\frac{e^{-\alpha u^2}}{d-iu}\,du
       =\pi e^{\alpha d^2}\operatorname{erfc}(d\sqrt\alpha).
\]
Indeed, \((d-iu)^{-1}=(d+iu)/(d^2+u^2)\), and the imaginary part
integrates to zero.

The small- and large-parameter limits are
\begin{equation}
 \frac{I(L)}{m_L}
       =b(0)+\frac{L}{8\pi}b''(0)+O(L^2)
                         \quad(L\downarrow0),\qquad
 \frac{I(L)}{m_L}=\frac12\log L+O(1)
                         \quad(L\to\infty).
 \label{v4-arith:weil:eq:auxiliary-limits}
\end{equation}
For the first, scale \(u=\sqrt{L/(2\pi)}v\) and use the second
moment \(1/2\) of \(\pi^{-1/2}e^{-v^2}\,dv\).
Taylor expansion near zero and the Gaussian tail justify the remainder.
For the second, use \(b(u)=\log|u|+O(1)\) outside a bounded interval.
After scaling, \(\log|v|\) is Gaussian-integrable.
The shrinking central interval contributes
\(O(L^{-1/2}\log L)\) after subtraction of \(\tfrac12\log L\).

To state the auxiliary cone in the \(2\pi\) convention, define
\[
 G_{2\pi}(f,g)=\frac1{2\pi}\int
       \widehat f_{2\pi}(u)\overline{\widehat g_{2\pi}(u)}b(u)\,du,
 \qquad F_L(x)=e^{-\pi Lx^2}.
\]
For \(L,M>0\), direct multiplication of their Fourier transforms gives
\begin{equation}
 G_{2\pi}(F_L,F_M)
       =\frac{I(2LM/(L+M))}{2\pi\sqrt{LM}}.
 \label{v4-arith:weil:eq:auxiliary-cross}
\end{equation}

\begin{proposition}[The exact auxiliary cone]
\label{v4-arith:weil:prop:auxiliary-cone}
The diagonal \(G_{2\pi}(F_L,F_L)\) is nonnegative exactly for
\(L\ge L_b\), and vanishes exactly at \(L=L_b\).
For a finite sum \(f=\sum_j a_jF_{L_j}\) with \(a_j\ge0\) and
\(L_j\ge L_b\), one has \(G_{2\pi}(f,f)\ge0\).
A nonzero such sum has zero value precisely when every active
parameter equals \(L_b\). Nonnegativity also holds on the
Schwartz closure of this cone.
\end{proposition}
\begin{proof}
The diagonal formula is \(I(L)/(2\pi L)\).
Its sign follows from \eqref{v4-arith:weil:eq:auxiliary-H}.
For \(L,M\ge L_b\), their harmonic mean is at least \(L_b\), so
every cross term in \eqref{v4-arith:weil:eq:auxiliary-cross}
is nonnegative. Expand the finite sum.

An active parameter strictly above \(L_b\) has a strictly positive
diagonal term. If all active parameters equal \(L_b\), every term
vanishes. Schwartz continuity proves the closure assertion.
\end{proof}

For example, take \(L_b/2<L<L_b\) and
\(M>L_bL/(2L-L_b)\). Then \(M>L_b\) and
\(2LM/(L+M)>L_b\), but \(I(L)<0\).
The value of \(G_{2\pi}(F_L+\varepsilon F_M,
F_L+\varepsilon F_M)\) remains negative for sufficiently small
\(\varepsilon>0\). This lies outside the stated cone.

Neither Gaussian cone is a complex vector space. Its nonzero elements
are positive on all of \(\mathbb R\), and hence are not compactly
supported. The span of centered Gaussians lies in the closed even
Schwartz subspace, so it is not dense in the full Schwartz space.
These facts do not prevent their use in \(\mathscr A_\infty\).
They specify which density and positivity claims are available.

For clarity, the two Fourier normalizations obey
\[
 \begin{aligned}
 \frac1{2\pi}\int\widehat f_{2\pi}(u)
       \overline{\widehat g_{2\pi}(u)}k(u)\,du
   &=\frac{\log\pi}{\pi}\langle f,g\rangle_{L^2}
                                      -G_{2\pi}(f,g),\\
 W_\infty(f*g^*)&=\int\widehat f_{2\pi}(u)
               \overline{\widehat g_{2\pi}(u)}k(2\pi u)\,du.
 \end{aligned}
\]
The second identity follows from \(h_f(2\pi u)=\widehat f_{2\pi}(-u)\).
A change of convention rescales the kernel argument as well as the
measure.

\section{The complete positivity criterion}
\label{v4-arith:weil:positivity}

The signs of the individual arithmetic terms do not decide the sign of
their sum. The spectral formula does decide it once every zero lies on
the critical line. The converse requires a test detecting a nonreal
ordinate while controlling every other zero.

Let \(\Sigma\) be the set of distinct ordinates
\(z_\rho=-i(\rho-1/2)\), and let \(m_z\) be their multiplicities.
The functional equation and conjugation make \(\Sigma\) invariant under
\(z\mapsto-z\) and \(z\mapsto\bar z\), preserving multiplicity.
It lies in \(|\operatorname{Im}z|\le1/2\) and is locally finite.
Equation~\eqref{v4-arith:weil:eq:pairing} becomes
\begin{equation}
 q_W(f,g)=\sum_{z\in\Sigma}m_z
                 h_f(z)\overline{h_g(\bar z)}.
 \label{v4-arith:weil:eq:distinct-pairing}
\end{equation}
This absolutely convergent formula holds on every exponential test step.

An individual nonreal pair has an indefinite two-dimensional form.
Polynomial interpolation and a Gaussian control the other coordinates.

\begin{theorem}[A negative test from an off-line zero]
\label{v4-arith:weil:thm:negative-test}
Suppose \(z_0=\gamma_0+i\eta_0\in\Sigma\), with \(\eta_0\ne0\).
Then a polynomial times a modulated Gaussian belongs to
\(\mathscr A_\infty\) and has negative Weil diagonal.
There is also a compact smooth test with negative Weil diagonal.
\end{theorem}
\begin{proof}
Choose \(R>1\), and let \(Z_R\) consist of the distinct points of
\(\Sigma\) with \(|\operatorname{Re}z-\gamma_0|\le R\).
It is finite and conjugation invariant. The polynomial
\[
 P_R(z)=\prod_{w\in Z_R\setminus\{z_0\}}
                  \frac{z-w}{z_0-w}
       -\prod_{w\in Z_R\setminus\{\bar z_0\}}
                  \frac{z-w}{\bar z_0-w}
\]
takes the values \(1,-1\) at \(z_0,\bar z_0\), respectively, and
vanishes on the other points of \(Z_R\). For \(\tau\ge1\), put
\begin{equation}
 \begin{aligned}
 H_\tau(z)&=P_R(z)
            e^{-\tau((z-\gamma_0)^2+\eta_0^2)},\\
 f_\tau(x)&=e^{-\tau\eta_0^2}P_R(i\partial_x)
         \frac{e^{-x^2/(4\tau)-i\gamma_0x}}{\sqrt{4\pi\tau}}.
 \end{aligned}
 \label{v4-arith:weil:eq:negative-test}
\end{equation}
Fourier differentiation gives \(h_{f_\tau}=H_\tau\).

Every derivative of \(f_\tau\) has Gaussian decay.
Its two selected transform values remain \(1,-1\), so their
contribution to the diagonal is \(-2m_{z_0}\).

For \(z=\gamma+i\eta\notin Z_R\), define
\[
 d(z)=(\gamma-\gamma_0)^2-\eta^2+\eta_0^2,
 \qquad c=R^2-\tfrac14+\eta_0^2>0.
\]
The strip bound gives \(d(z)\ge c\). The absolute value of the
remaining contribution is at most
\begin{equation}
 e^{-2(\tau-1)c}
   \sum_{z\notin Z_R}m_z|P_R(z)P_R(\bar z)|e^{-2d(z)}.
 \label{v4-arith:weil:eq:negative-tail}
\end{equation}
The sum is finite.

Indeed, the polynomial has polynomial growth
on the fixed strip, the zero count is \(O(T\log(2+T))\), and the
Gaussian controls the resulting sum over unit intervals.
Choose \(\tau\) large enough that \eqref{v4-arith:weil:eq:negative-tail}
is less than \(m_{z_0}\). Then \(q_W(f_\tau,f_\tau)<-m_{z_0}\).

Fix this test and multiply it by compact smooth cutoffs tending to one.
The products converge in each \(\mathscr A_{a_0}\).
The continuity in Theorem~\ref{v4-arith:weil:thm:extension} preserves
strict negativity for all sufficiently large cutoffs.
\end{proof}

\begin{theorem}[Weil positivity on its full domain]
\label{v4-arith:weil:thm:criterion}
The following assertions are equivalent:
\begin{enumerate}
 \item Every zero of \(\xi\) has real part \(1/2\).
 \item \(q_W(f,f)\ge0\) for every \(f\in\mathscr A_c\).
 \item \(q_W(f,f)\ge0\) for every \(f\in\mathscr A_{\exp}\).
 \item For some fixed \(a_0>1/2\), a subset dense in
       \(\mathscr A_{a_0}\) has nonnegative Weil diagonal.
\end{enumerate}
The second assertion is the complete arithmetic inequality
\begin{equation}
 W_\infty(f*f^*)\le P(f*f^*)-W_{\mathrm{fin}}(f*f^*)
                \qquad(f\in\mathscr A_c).
 \label{v4-arith:weil:eq:complete-inequality}
\end{equation}
\end{theorem}
\begin{proof}
If all ordinates are real, \eqref{v4-arith:weil:eq:distinct-pairing}
is a sum of absolute squares on every exponential step.
This implies the third assertion and hence the second and fourth.
Continuity extends the fourth assertion to its whole step, which
contains all compact smooth tests. Theorem~\ref{v4-arith:weil:thm:negative-test}
shows that the second assertion excludes every nonreal ordinate.
Finally, the arithmetic definition of \(W\) gives the displayed inequality.
\end{proof}

\begin{proposition}[Hilbert realization of a positive restriction]
\label{v4-arith:weil:prop:hilbert}
Let \(D\subset\mathscr A_{\exp}\) be a complex vector subspace on
which \(q_W\) is positive semidefinite. Its nullspace
\(N_D=\{f\in D:q_W(f,f)=0\}\) is a linear subspace.
The completion of \(D/N_D\) is a Hilbert space \(H_D\), and the
canonical linear map \(I_D:D\to H_D\) satisfies
\[
 \langle I_Df,I_Dg\rangle=q_W(f,g).
\]
Conversely, any such map to a positive Hilbert space implies positivity
on \(D\). In particular, existence of this realization on the full
compact test algebra is equivalent to the Riemann hypothesis.
\end{proposition}
\begin{proof}
Positivity of \(q_W(f+tg,f+tg)\) for complex \(t\) gives
Cauchy--Schwarz. A null vector therefore pairs to zero with every
vector. This proves that the quotient form is well-defined and positive
definite. Its metric completion has the asserted property.
The converse follows by taking the diagonal, and the last assertion
uses Theorem~\ref{v4-arith:weil:thm:criterion}.
\end{proof}

This quotient construction starts from the Weil form. It supplies no
independent arithmetic reason for its positivity. An arithmetic Hilbert
space, constructed by a different method, must still be connected to it
by the displayed pairing identity on the stated domain.

The compact domain also distinguishes this criterion from rational
tests used in Li's criterion. For an integer \(n\ge1\),
\[
 1-(1-1/s)^n=\frac ns+O(s^{-2})\qquad(|s|\to\infty).
\]
After \(s=1/2+iz\), this is neither entire nor rapidly decreasing
on horizontal strips. It is not \(h_f\) for a compact smooth test.
A formula involving these rational tests needs its own limiting
prescription. It is not obtained by substituting them into the compact
explicit formula without a convergence argument.

\section{Moment equations and the topology of approximation}
\label{v4-arith:weil:moments}

A proposed arithmetic comparison may first be defined on tests satisfying
moment equations. Their closed kernel can often be approximated exactly,
but this does not make it dense in the full domain. We give the projection
and determine its closure in two different topologies.

For a compact test put
\[
 M_f(s)=\int_{\mathbb R}f(x)e^{sx}\,dx=h_f(-is).
\]
On the real even compact tests, define
\[
 \ell(f)=\left(\int f(x)\,dx,\
          \int f(x)\cosh x\,dx,\
          \int f(x)x\sinh x\,dx\right),\qquad K=\ker\ell.
\]
Thus \(K\) imposes \(M_f(0)=M_f(1)=M_f'(1)=0\).
Evenness already gives \(M_f'(0)=0\).
These are the double-zero equations at \(0,1\) for an even transform.

All three functionals are continuous on \(\mathscr A_\infty\).
They are also continuous on each \(\mathscr A_{a_0}\) with \(a_0>1\).

\begin{proposition}[A continuous moment projection]
\label{v4-arith:weil:prop:moment-projection}
There are real even compact tests \(\eta_0,\eta_1,\eta_2\) with
\(\ell_i(\eta_j)=\delta_{ij}\). Hence
\begin{equation}
 \Pi f=f-\sum_{j=0}^2\ell_j(f)\eta_j
 \label{v4-arith:weil:eq:moment-projection}
\end{equation}
is a continuous projection onto the moment kernel. That kernel is closed
of codimension three, both on compact tests and on each of the exponential
spaces just specified. For any \(A>0\), the \(\eta_j\) may all be chosen
with support in \([-A,A]\).
\end{proposition}
\begin{proof}
Write \(w(x)=(1,\cosh x,x\sinh x)\).
The matrix with columns \(w(0),w(1),w(2)\) has determinant
\[
 2\sinh(1)(\cosh(1)-1)^2>0.
\]
Even mollified point masses at \(0,\pm1,\pm2\), with total mass
one at each pair, give a moment matrix converging to this matrix.
For small mollifier support, its inverse supplies the \(\eta_j\).
Continuity of the three moments then proves the projection and
codimension assertions.

To prescribe the support, choose a smooth even \(\chi\ge0\), supported
in \([-A,A]\) and positive on \((-A,A)\). The Gram matrix
\[
 G_{ij}=\int\chi(x)w_i(x)w_j(x)\,dx
\]
is positive definite. A nonzero linear combination of the three analytic
functions cannot vanish on an interval. Analytic continuation and their
distinct constant, exponential and linear-exponential growth show their
independence. Set
\(\eta_j=\chi\sum_i(G^{-1})_{ij}w_i\).
These tests have the required moments and support.
\end{proof}

If \(f_j\to f\in K\) in the compact-test topology, then
\(\Pi f_j\to f\). In an exponential space, even compact cutoffs followed
by \(\Pi\) prove that compact members of the kernel are dense in that
kernel. A nonnegative bump has nonzero mean, so the kernel is not dense
in the ambient compact-test space. A moment correction converges to zero
only when the limiting function already has the required moments.

\begin{proposition}[The finite moment quotient]
\label{v4-arith:weil:prop:moment-quotient}
On the complex even compact convolution algebra the map
\begin{equation}
 f\longmapsto\bigl(M_f(0),M_f(1)+M_f'(1)\varepsilon\bigr),
 \qquad \varepsilon^2=0,
 \label{v4-arith:weil:eq:moment-quotient}
\end{equation}
is a continuous surjective algebra homomorphism onto
\(\mathbb C\times\mathbb C[\varepsilon]/(\varepsilon^2)\).
Its kernel is the complexification of \(K\), and its involution is
coefficient conjugation with \(\varepsilon^*=\varepsilon\).
A complex-linear functional on this quotient is positive on squares
exactly when it has the form
\[
 \lambda(a,b+c\varepsilon)=\alpha a+\beta b,
 \qquad \alpha,\beta\ge0.
\]
\end{proposition}
\begin{proof}
Fubini gives \(M_{f*g}=M_fM_g\). Evaluation at zero and the product
rule at one give the algebra law. The moment right inverse proves
surjectivity. Evenness makes the involution conjugation at these real
evaluation points.

Write \(\lambda(a,b+c\varepsilon)=\alpha a+\beta b+\gamma c\).
Positivity on the unital quotient makes \(\lambda\) Hermitian, by
expanding squares of \(1+t x\) and \(1+itx\) for real \(t\).
Thus its coefficients are real. Its value at the square of
\((a,b+c\varepsilon)\) is
\[
 \alpha|a|^2+\beta|b|^2+2\gamma\operatorname{Re}(b\bar c).
\]
Taking \(b=1\) and varying \(c\) forces \(\gamma=0\).
The two remaining coefficients must be nonnegative, which also suffices.
\end{proof}

This quotient describes the imposed moments. It does not assert that
\(W\) factors through the quotient. The polar evaluations occur at
\(\pm1/2\), not at \(0,1\). For a nonzero nonnegative even compact
test \(v\), the function
\[
 f=D^2(D^2-1)^2v,\qquad D=d/dx,
\]
lies in \(K\), whereas
\[
 M_f(1/2)=\frac9{64}M_v(1/2)>0,
 \qquad P(f*f^*)=2M_f(1/2)^2>0.
\]
Integration by parts gives
\(M_f(s)=s^2(s^2-1)^2M_v(s)\), proving every assertion.

\begin{theorem}[The weaker Schwartz closure]
\label{v4-arith:weil:thm:moment-closure}
The closure of \(K\) in the real even Schwartz space is
\[
 \left\{f\in\mathcal S(\mathbb R):f\text{ is real and even},\
                                    \int f=0\right\}.
\]
In the relative Schwartz topology on compact tests, intersect this set
with \(\mathscr A_c\). The two exponential moment equations need not
survive a limit in this weaker topology.
\end{theorem}
\begin{proof}
The mean is Schwartz-continuous, so it vanishes on every such limit.
Choose a nonnegative even compact test \(\beta\) with \(\int\beta=1\),
and put
\[
 A=\int\beta(x)\cosh x\,dx>0,\qquad
 B=\int\beta(x)x\sinh x\,dx.
\]
For \(R>0\), set
\[
 \beta_R(x)=\tfrac12[\beta(x-R)+\beta(x+R)],\qquad
 v_R=\beta_R-\beta.
\]
Translation gives
\[
 \ell(v_R)=\bigl(0,A(\cosh R-1),
              B(\cosh R-1)+AR\sinh R\bigr).
\]

The determinant of the last two coordinates of \(v_R,v_{2R}\) is
\[
 2A^2R\sinh R(\cosh R-1)^2>0.
\]
Thus a fixed pair of these moments can be prescribed by
\(c_Rv_R+d_Rv_{2R}\), where
\(c_R=O(e^{-R})\) and \(d_R=O(e^{-2R})\).
Every Schwartz seminorm of \(v_R\) grows at most polynomially in \(R\).
The correction therefore tends to zero in the Schwartz topology.

For a compact real even mean-zero \(f\), prescribe the negatives of its
last two moments. Adding the correction gives elements of \(K\)
converging to \(f\). General real even mean-zero Schwartz functions are
approximated by even compact cutoffs with their means subtracted times
\(\beta\). For the \(n\)th approximation, choose \(R\) so large that
the correcting term has its first \(n\) Schwartz seminorms below \(1/n\).
This diagonal choice proves the asserted closure.
\end{proof}

The complete Weil form is continuous in the compact and exponential
topologies used above. The critical prime measure is not tempered by
Theorem~\ref{v4-arith:weil:thm:prime-threshold}. Consequently, the weaker
Schwartz closure does not itself pass the complete Weil inequality to
its limits.

\begin{proposition}[The missing cross inequalities]
\label{v4-arith:weil:prop:cross}
Let a complex vector space decompose as \(D=K_0\oplus E\), and let
\(q\) be a Hermitian form, linear in its first argument. It is positive
semidefinite on \(D\) exactly when
\[
 q(u,u)\ge0,\qquad q(e,e)\ge0,\qquad
 |q(u,e)|^2\le q(u,u)q(e,e)
 \quad(u\in K_0,\ e\in E).
\]
\end{proposition}
\begin{proof}
Necessity is Cauchy--Schwarz. Conversely, the three conditions give
\[
 q(u+e,u+e)\ge
       \bigl(\sqrt{q(u,u)}-\sqrt{q(e,e)}\bigr)^2\ge0.
\]
Every vector has this form.
\end{proof}

The complex moment kernel has a three-dimensional complement among even
tests. Its positivity leaves both that complement and all cross terms
to determine. Odd tests remain an additional part of the full domain.
The exact polar kernel has a two-dimensional complement on the full
space. Finite codimension alone supplies none of these inequalities.

Similarly, if maps \(I_\alpha:D_\alpha\to H_\alpha\) to positive
Hilbert spaces satisfy
\(q_W(f,g)=\langle I_\alpha f,I_\alpha g\rangle\) on their respective
complex domains, their union has nonnegative diagonal. Its closure in a
fixed exponential step has the same property. Their linear span requires
the cross terms between different domains. The Hermitian matrix
\(\left(\begin{smallmatrix}1&2\\2&1\end{smallmatrix}\right)\) is positive
on each coordinate axis and negative on \((1,-1)\). Thus separate
positive diagonal identities do not suffice for the span.

\section{The maximal positive convolution ideal}
\label{v4-arith:weil:positive-ideal}

Convolution stability gives a stronger classification than positivity
on a subspace. A positive ideal must vanish at every nonreal ordinate.
To prove this assertion, we first construct tests with exactly prescribed
values on a finite part of the zero set.

\begin{lemma}[A paired product and its Fourier inverse]
\label{v4-arith:weil:lem:paired-product}
The product
\begin{equation}
 \mathcal P(z)=
  \prod_{\substack{\omega\in\Sigma\\\operatorname{Re}\omega>0}}
             (1-z^2/\omega^2)^{m_\omega}
 \label{v4-arith:weil:eq:paired-product}
\end{equation}
converges locally uniformly to an entire function with precisely the
zeros \(\Sigma\) and their multiplicities. It is real on the real axis
and satisfies
\[
 |\mathcal P(z)|\le
       \exp\bigl(C(1+|z|)\log(2+|z|)\bigr).
\]
For a finite \(E\subset\Sigma\), division by
\(\prod_{\omega\in E}(z-\omega)^{m_\omega}\) gives an entire function
\(R_E\), nonzero at each point of \(E\), with the same type of bound.
For \(t>0\) and any polynomial \(p\),
\(e^{-tz^2}R_E(z)p(z)\) is the Fourier transform of a member of
\(\mathscr A_\infty\).
\end{lemma}
\begin{proof}
There are no ordinates on the imaginary axis. Pairing \(\omega\) and
\(-\omega\) therefore accounts for all zeros. Inverse-square summability
proves normal convergence of the quadratic factors. Conjugation permutes
these factors and proves reality on the real axis.

Let \(n(u)\) count the factors by modulus with multiplicity, and choose
\(u_0>0\) below the modulus of every zero. For \(R\ge2\),
\[
 \begin{aligned}
 \log\max_{|z|\le R}|\mathcal P(z)|
 &\le\int_{u_0}^\infty\log(1+R^2/u^2)\,dn(u)\\
 &=2R^2\int_{u_0}^\infty
                    \frac{n(u)}{u(u^2+R^2)}\,du.
 \end{aligned}
\]
Split at \(R\) and use \(n(u)=O(u\log(2+u))\).
Both integrals are \(O(R\log R)\).

Removing the full order of each selected zero leaves a nonzero value
there. Outside a fixed disk the dividing polynomial has a polynomial
lower bound. Inside that disk the entire quotient is bounded.

Put \(H(z)=e^{-tz^2}R_E(z)p(z)\). On each horizontal strip, the negative
quadratic in \(\operatorname{Re}z\) dominates the other factors.
Cauchy's estimates on a slightly larger strip give rapid decay of
all derivatives. Define
\[
 f(x)=\frac1{2\pi}\int_{\mathbb R}H(r)e^{-irx}\,dr.
\]
For \(x>0\) shift to \(\operatorname{Im}z=-b\), and for \(x<0\)
to \(\operatorname{Im}z=b\). The vertical integrals tend to zero.

The same shift applied to \(z^jH(z)\) gives
\(|f^{(j)}(x)|\le C_{b,j}e^{-b|x|}\).
Since \(b\) is arbitrary, \(f\in\mathscr A_\infty\).
Fourier inversion and holomorphic continuation give \(h_f=H\).
\end{proof}

\begin{theorem}[Exact finite-value isolation]
\label{v4-arith:weil:thm:isolation}
For a finite \(E\subset\Sigma\) and prescribed complex numbers
\(v_\omega\), there is \(f\in\mathscr A_\infty\) such that
\(h_f(\omega)=v_\omega\) on \(E\) and \(h_f=0\) on
\(\Sigma\setminus E\).
\end{theorem}
\begin{proof}
For any \(t>0\), take the inverse Fourier transform of
\begin{equation}
 e^{-tz^2}R_E(z)
 \sum_{\omega\in E}\frac{v_\omega e^{t\omega^2}}{R_E(\omega)}
        \prod_{\substack{\eta\in E\\\eta\ne\omega}}
                          \frac{z-\eta}{\omega-\eta}.
 \label{v4-arith:weil:eq:isolation}
\end{equation}
All denominators are nonzero. The polynomial interpolates the required
values after division by the nonvanishing prefactor, and \(R_E\) retains
every other zero. Lemma~\ref{v4-arith:weil:lem:paired-product} supplies
the required inverse transform.
\end{proof}

Full multiplicities are removed in \(R_E\), although the prescription
specifies only values. Removing one linear factor at a multiple zero
would leave the prefactor zero there. Exact isolation here uses rapid
exponential decay. Compact cutoffs preserve a strict inequality, but
need not preserve these exact spectral values.

Write \(\Sigma_0=\Sigma\cap\mathbb R\) and
\(\Sigma_{\mathrm{nr}}=\Sigma\setminus\mathbb R\).
A positive convolution \(*\)-ideal means a complex algebraic ideal
\(J\subset\mathscr A_{\exp}\), stable under \(*\), with
\(q_W(f,f)\ge0\) for every \(f\in J\).
Define
\[
 \mathcal I_{\mathrm{nr}}=
 \{f\in\mathscr A_{\exp}:h_f(z)=0
                         \text{ for every }z\in\Sigma_{\mathrm{nr}}\}.
\]

\begin{theorem}[The largest positive ideal]
\label{v4-arith:weil:thm:maximal-ideal}
The ideal \(\mathcal I_{\mathrm{nr}}\) is the largest positive convolution
\(*\)-ideal of \(\mathscr A_{\exp}\). It is closed in the final test
topology and in every exponential step. On this ideal,
\begin{equation}
 W(f)=\sum_{\gamma\in\Sigma_0}m_\gamma h_f(\gamma),\qquad
 q_W(f,g)=\sum_{\gamma\in\Sigma_0}m_\gamma
                         h_f(\gamma)\overline{h_g(\gamma)}.
 \label{v4-arith:weil:eq:positive-ideal}
\end{equation}
If a \(*\)-ideal contains \(f\) with \(h_f(z_0)\ne0\) at a nonreal
ordinate, it contains \(r\) with \(q_W(r,r)=-2m_{z_0}\).
\end{theorem}
\begin{proof}
Each evaluation at a point of \(\Sigma\) is continuous on every step,
hence on the final topology. Fourier multiplication and conjugation
show that the intersection of the indicated kernels is a closed
\(*\)-ideal. The explicit formula gives
\eqref{v4-arith:weil:eq:positive-ideal} and its positivity.

For maximality put \(u=h_f(z_0)\ne0\). Exact isolation supplies
\(e_{z_0},e_{\bar z_0}\) whose transforms on \(\Sigma\) are the respective
coordinate vectors. In the given ideal form
\[
 r=u^{-1}f*e_{z_0}-\bar u^{-1}f^**e_{\bar z_0}.
\]
Its only nonzero transform values are \(1\) at \(z_0\) and \(-1\)
at \(\bar z_0\), since \(h_{f^*}(\bar z_0)=\bar u\).
Their equal multiplicities give the asserted negative value.
Every positive ideal must therefore lie in \(\mathcal I_{\mathrm{nr}}\).
\end{proof}

This theorem classifies positive ideals. A positive subspace need not
contain the two convolution products used in its proof, so it need not
lie in \(\mathcal I_{\mathrm{nr}}\). The theorem neither asserts nor
excludes the existence of a nonreal zero.

The regular completions lose the distinction between this ideal and the
whole test algebra. To see this, form the off-real subproduct
\[
 B_{\mathrm{nr}}(z)=
 \prod_{\substack{\omega\in\Sigma\\
          \operatorname{Re}\omega>0,\ \operatorname{Im}\omega>0}}
       \bigl[(1-z^2/\omega^2)(1-z^2/\bar\omega^2)\bigr]^{m_\omega}.
\]
The empty product equals one. The preceding product estimate applies.

Its zeros are precisely the nonreal ordinates with their multiplicities,
and for real \(r\),
\[
 B_{\mathrm{nr}}(r)=
 \prod_{\substack{\operatorname{Re}\omega>0\\
                   \operatorname{Im}\omega>0}}
              |1-r^2/\omega^2|^{2m_\omega}>0.
\]
For fixed \(t>0\), let \(b_t\in\mathscr A_\infty\) have transform
\(\beta_t(z)=e^{-tz^2}B_{\mathrm{nr}}(z)\), and put
\(\mathcal D_t=b_t*\mathscr A_{\exp}\).
This is a principal ideal in an algebra without a unit.

\begin{theorem}[Density in the regular completions]
\label{v4-arith:weil:thm:regular-density}
The ideal \(\mathcal D_t\) is a positive convolution \(*\)-ideal
contained in \(\mathcal I_{\mathrm{nr}}\). Its Fourier image is dense
in \(\mathcal S(\mathbb R)\). Consequently,
\[
 \overline{\mathcal D_t}^{L^1}=L^1(\mathbb R),\qquad
 \overline{\mathcal D_t}^{L^2}=L^2(\mathbb R),\qquad
 \overline{h(\mathcal D_t)}^{\|\cdot\|_\infty}=C_0(\mathbb R).
\]
The same three closures hold for \(\mathcal I_{\mathrm{nr}}\).
If a nonreal ordinate exists, \(\mathcal I_{\mathrm{nr}}\) is nevertheless
proper and closed in the test topology.
\end{theorem}
\begin{proof}
The function \(\beta_t\) is real on the real axis, so \(b_t^*=b_t\).
Commutativity makes \(\mathcal D_t\) a \(*\)-ideal. Its Fourier factor
vanishes on every nonreal ordinate, proving inclusion and positivity.

Take \(H\in C_c^\infty(\mathbb R)\) in the Fourier variable, and put
\(u=H/\beta_t\). Strict positivity of \(\beta_t\) on the real axis
makes \(u\) smooth and compactly supported. Define
\[
 \varphi_\epsilon(r)=(4\pi\epsilon)^{-1/2}e^{-r^2/(4\epsilon)},
 \qquad u_\epsilon=u*\varphi_\epsilon,
 \qquad H_\epsilon=\beta_tu_\epsilon.
\]
Convolution here is in the Fourier variable. Its inverse transform is
\(\check u_\epsilon(x)=e^{-\epsilon x^2}\check u(x)\), which belongs
to \(\mathscr A_\infty\). Thus \(H_\epsilon\in h(\mathcal D_t)\).

Gaussian convolution converges to the identity in the Schwartz topology.
This follows from translation continuity near zero and the polynomially
weighted Gaussian tail outside a fixed neighborhood. Multiplication by
the fixed Schwartz function \(\beta_t\) is continuous, so
\(H_\epsilon\to H\) in that topology.
Smooth cutoffs make compact smooth functions dense in Schwartz space.
Fourier inversion and its continuous dense inclusions into \(L^1\),
\(L^2\) and \(C_0\) give the three closures.

If \(z_0\) is nonreal, the continuous evaluation \(h_f(z_0)\) vanishes
on \(\mathcal I_{\mathrm{nr}}\) but not on a Gaussian. This proves
properness in the test topology.
\end{proof}

Neither ideal has an internal unit. Such a unit would act identically
on a uniformly dense Fourier subalgebra of \(C_0(\mathbb R)\).
Its transform would then equal one on all of \(\mathbb R\).
Nor is test-topology density of \(\mathcal D_t\) in
\(\mathcal I_{\mathrm{nr}}\) asserted. The generator vanishes to full
analytic multiplicity, while the latter ideal requires only value
vanishing.

These ideals depend on the zero set. Their unconditional construction
does not supply an arithmetic positive domain independently of that set.
More decisively, density in any one exponential test step would imply
full positivity by Theorem~\ref{v4-arith:weil:thm:criterion}.
The regular density theorem therefore proves a different assertion.

\section{The independent arithmetic comparison}
\label{v4-arith:weil:comparison}

The analytic target is now fixed: the convolution algebra, its topology,
its involution, and the complete signed functional. An arithmetic source
may approach this target through a trace, a cohomological pairing, an
automorphic norm or a boundary construction. Each approach must specify
its map on the same tests.

\begin{question}[A full arithmetic realization]
\label{v4-arith:weil:qn:comparison}
Construct an arithmetic complex vector space \(V\), a Hermitian pairing
\(B\), and a complex-linear map \(I:\mathscr A_c\to V\), such that
\begin{equation}
 B(If,Ig)=P(f*g^*)-W_\infty(f*g^*)-W_{\mathrm{fin}}(f*g^*)
                   \qquad(f,g\in\mathscr A_c).
 \label{v4-arith:weil:eq:arithmetic-comparison}
\end{equation}
Determine whether \(B\) is positive semidefinite on \(I(\mathscr A_c)\)
by the arithmetic construction. For an exponential extension, construct
the continuous map and prove the same identity on a specified step.
\end{question}

An operator or weight realization needs an actual convolution
representation on a common domain that preserves the involution.
Its trace or weight must be finite on the tested products. A comparison
with a bilinear categorical pairing also needs the conjugate coefficient
factor that turns it into a Hermitian form. Local determinants, a scalar
functional equation and matching characters do not specify these maps.

A geometric realization asks for a continuous \(*\)-algebra map into a
proper positive trace ideal constructed from independent arithmetic
operations. The comparison must preserve the involution, normal
derivatives, collision maps and both vacuum data of its source.
Arithmetic-weight and arity-radius continuity are additional requirements.
The resulting counting weight or self-adjoint shift must be identified
with a native arithmetic operation. Positivity of the scalar Weil form
alone does not construct this stronger realization.

Partial domains retain their exact positive content. A norm identity on
a moment kernel proves positivity on that kernel. A Gaussian cone gives
its tested inequalities. A positive ideal has the classification above.
To reach the full criterion, one must also prove the missing cross
inequalities and density in a topology controlling every term of
\eqref{v4-arith:weil:eq:arithmetic-comparison}.
Both polar evaluations, all prime powers and the signed gamma integral
remain part of that equality.
